% Supplementary Notes — resolved review/audit items (synced to P0 round-17 state, 2026-09-09)
\appendix
\section{Supplementary Notes}

\subsection{Supplementary Note: QC-metrics-only baseline}
\label{sn:qc_baseline}

To test whether predictions could be explained by library-quality or
processing artifacts rather than biology, we trained logistic regression on
per-patient QC metrics (mean UMI counts, detected genes, percent mitochondrial
reads, and cell count) under the identical repeated-CV protocol
(Section~\ref{ssec:methods_cv}). On the pooled analysis cohort (252
patients, 27 fatal), the QC-metrics-only model reached a patient-level AUC of
0.528 $\pm$ 0.026 (repeat-mean $\pm$ repeat-sd);
on the COVID-only cohort (240 patients, 25 fatal) it scored below chance
(0.376 as-fit). The below-chance value is a textbook suppression pattern, not
an anomaly: the QC features are mechanically correlated with library depth
and carry outcome information only conditional on it. Adding each patient's
log-transformed cell count to the same model recovers a correctly oriented
0.618 as-fit, and adding cohort identity on top reaches 0.627; it is the cell
count, not cohort, that flips the sign (cohort contributes $+0.009$). Fatal
patients showed higher mean UMI counts and mitochondrial fractions than
survivors in both diseases, consistent with more activated, fragment-rich
platelet transcriptomes (Supplementary Table~\ref{tab:qc_strat}). Because a
baseline race cannot separate these contributions, the confound question is
settled by the incremental-value test in Supplementary
Note~\ref{sn:qc_incremental}.

\subsection{Supplementary Note: incremental value over QC, cell-count, and cohort covariates}
\label{sn:qc_incremental}

The decisive confound test is incremental value, not a baseline race. The
baseline covariate set (per-patient mean UMI counts, detected genes, percent
mitochondrial reads, log-transformed cell count, and cohort indicator
variables; 13 covariates) was fit with L2-regularized logistic regression
($C = 0.3$; regularization is required at 13 covariates and 25 events) under
the identical repeated-CV protocol, reaching a patient-level AUC of 0.627
as-fit on the seed-ensemble scores.

The pre-specified primary test for this question is a single nested
likelihood-ratio test, covariates-only against covariates plus the PCA-200
representation score, chosen because it involves no model selection and is
exactly the question asked: does a platelet representation add over QC, cell
count, and cohort? Because an L2-penalized statistic has no valid $\chi^2(1)$
reference, the null distribution is empirical: patient-level label
permutation (10{,}000 permutations) with both nested models refit per
permutation at the identical penalty. The observed statistic (4.477) exceeds
every permuted maximum, $P < 0.0001$ at the estimator's resolution floor
($P < 1/10{,}001$). This pre-specified test alone carries the conclusion that
mortality is learnable from platelet transcriptomes beyond QC, cell count,
and cohort.

The five foundation models are reported as exploratory, with two readings
paired so that neither is quoted alone (Supplementary
Table~\ref{tab:incremental_value}). By the same nested likelihood-ratio test
against the empirical null, every representation added significant value
(four of five at the resolution floor $P < 0.0001$; \model{UCE} weakest at
$P = 0.0022$); by paired-AUC comparison against the same baseline, five of
six increments were significant (\model{STATE} $+0.065$, \model{Geneformer}
$+0.053$, PCA-200 $+0.050$, \model{scPRINT-2} $+0.046$, \model{scGPT}
$+0.026$; exact DeLong $p = 0.0001$--$0.006$), with \model{UCE} again the
exception ($+0.016$, $p = 0.072$). The likelihood-ratio family establishes
that incremental information exists; the AUC increments quantify its
magnitude; the penalized statistic is sensitive to presence but nearly
insensitive to magnitude, so no effect size should be read from the
likelihood-ratio $P$-values themselves. Two diagnostic tables accompany the
result: the cross-validation folds are cohort-balanced (Supplementary
Table~\ref{tab:fold_balance}), and the per-cohort QC--mortality associations
are reported in Supplementary Table~\ref{tab:qc_by_cohort}.

\subsection{Supplementary Note: cohort-only baseline}
\label{sn:cohort_baseline}

Because cohort is the variable the QC diagnostic points at, we scored each
patient by their cohort's mortality rate, estimated within training folds
only with add-one (Laplace) smoothing, under the identical repeated-CV
protocol. The cohort-only model reached a patient-level AUC of
0.486 $\pm$ 0.038 (repeat-mean $\pm$ repeat-sd), statistically
indistinguishable from chance. This number must not be read as exonerating
cohort: add-one smoothing is size-dependent, so a small zero-death cohort
receives a higher smoothed rate than a large cohort with real deaths,
inverting the risk ordering by construction (per-cohort counts and smoothed
rates in Supplementary Table~\ref{tab:cohort_rates}). The estimate therefore
measures the failure of one particular cohort-risk estimator and is silent on
the confound question; the evidence on that question is the incremental-value
test with cohort identity in the covariate set
(Supplementary Note~\ref{sn:qc_incremental}).

\subsection{Supplementary Note: nested-comparison (complementarity) family}
\label{sn:nested_comparison}

A question distinct from superiority is complementarity: whether a
foundation-model score adds information beyond the baseline's score even
without ranking patients better on its own. We tested it with nested models:
a reduced model containing the PCA-200 score (in three comparator
configurations: the PCA-200 score alone; plus log-transformed cell count;
plus the full QC and cohort covariate set) against a full model adding each
foundation-model score in turn, refit inside the repeated-CV protocol with
all preprocessing fit within training folds, including per-fold feature
standardization. Any nested out-of-fold pooling harness of this kind is
required to pass a unit test before use: a known frozen score refit through
the harness must reproduce its own out-of-fold AUC within 0.02--0.04, and a
harness failing the gate produces no significance conclusions. With the gated
harness, the reduced models reached 0.710--0.724 patient-level AUC against
PCA-200's own 0.725, and no foundation model's increment survived
Benjamini--Hochberg correction ($\alpha = 0.05$) in any of the three
comparator configurations (0/5 in each; rank-by-rank values in Supplementary
Table~\ref{tab:nested_oof}); the only nominal signal was \model{STATE} at
$p = 0.025$--$0.045$ across configurations, not surviving correction. A
winner-corrected deviance permutation test on the same nested construction is
retained as an in-sample cross-check (Supplementary
Table~\ref{tab:nested_deviance}); because it does not admit an out-of-fold
form, it carries no headline weight.

\subsection{Supplementary Note: each method against the count-plus-disease null}
\label{sn:method_vs_null}

Each method's seed-ensemble scores were compared with the null model
(per-patient cell count and disease label; AUC 0.554) by exact DeLong paired
tests. Five of six methods significantly outscored the null: \model{STATE}
($p = 0.0014$), \model{Geneformer} ($p = 0.0006$), \model{scPRINT-2}
($p = 0.011$), \model{scGPT} ($p = 0.023$), and PCA-200 ($p = 0.013$), all
surviving Benjamini--Hochberg correction across the six tests. \model{UCE}
did not ($+0.122$, $p = 0.157$) and is the only method failing that
correction; it is the weakest method throughout this benchmark
(Section~\ref{ssec:r1_within_covid}).

\subsection{Supplementary Note: confirmatory range-statistic test for the leaders--UCE gap}
\label{sn:range_test}

The single clearly resolvable gap in the within-COVID ranking, the two
leaders' advantage over \model{UCE} ($+0.118$ and $+0.130$), was tested
against the winner-selection null with a range statistic on the same 10{,}000
patient-level label permutations as the primary analysis: the max-minus-min
AUC across the five foundation models was recorded per permutation, and the
observed range (0.130) exceeded the null distribution
($P = 0.011$). We report this separation as the paper's single confirmatory
secondary result (Section~\ref{ssec:r1_within_covid}).

\subsection{Supplementary Note: winner-corrected primary test}
\label{sn:winner_test}

The pre-specified primary question, whether the best foundation model beats
the from-scratch baseline, was tested with a max-statistic permutation test on
the seed-ensemble out-of-fold scores: mortality labels were permuted at the
patient level (10{,}000 permutations), the maximum
foundation-model-minus-baseline AUC difference was recorded for each
permutation, and the familywise $P$ is the fraction of permutations whose
maximum reaches the observed maximum. This construction simulates the null
while preserving the inter-model score correlations (Pearson $r = 0.71$--$0.90$;
Supplementary Table~\ref{tab:score_correlation}) and the winner-selection
procedure itself; because labels are permuted on frozen scores without
retraining, it tests the slightly weaker null that the scores carry no
outcome information. The observed maximum contrast (\model{Geneformer},
$+0.064$ AUC) gave $P = 0.181$ (Table~\ref{tab:winner_ensemble}); a
recentered-bootstrap cross-check (subtracting the observed per-model
differences before taking the maximum in each patient resample) gave
$P = 0.073$. The pre-specified primary contrast is therefore not met. As an
exploratory comparison, a consensus predictor averaging the five models'
scores reached AUC 0.768, above the PCA-200 baseline (0.742) but below the
best single model (\model{Geneformer} 0.807): averaging over the weaker
candidates scores 0.039 AUC below the post-hoc winner on the same patients
used to identify it; this comparison is descriptive, and a prospective choice
between winner selection and ensembling would require nested evaluation.

\subsection{Supplementary Note: stability of the primary result}
\label{sn:fragility}

Leave-one-fatal-patient-out refits (25 in total) repeated the permutation
max-statistic test on 239 patients. Every refit returned $P > 0.05$ (range
0.141--0.288): the primary null is not driven by any single death
(Section~\ref{ssec:r1_within_covid}).

\subsection{Supplementary Note: variance-estimation consistency check}
\label{sn:se_check}

Pairwise AUC contrasts are reported with exact DeLong standard errors
(placement-value covariance) computed on the seed-ensemble scores. As a
consistency check, patient-level paired bootstrap standard errors agree with
the exact DeLong values for all 15 within-COVID pairs (within about 5\%;
Supplementary Table~\ref{tab:se_check}). The Hanley--McNeil formula with the
raw-score Pearson correlation is an unpaired approximation and overestimates
the paired standard errors by 20--35\% in 8 of 15 pairs; it is not used for
any reported contrast. The winner-selection and range statistics
(Supplementary Notes~\ref{sn:winner_test} and~\ref{sn:range_test}) do not
admit a DeLong form and use permutation nulls throughout.

\subsection{Supplementary Note: age structure of the GSE279452 sepsis cohort}
\label{sn:age}

The analysis set and the GEO series use three different counts, which we
state once to prevent misreading: 93 is the number of Adult Sepsis candidate
donors in GSE279452; 92 of them matched per-sample age records directly (one
donor, \texttt{Skin-S152}, matches only after correcting a sample-title
spelling inconsistency in the GEO series, where the sample is titled
\texttt{Skin-152} while its same-site siblings carry the \texttt{Skin-S}
prefix); and 91 patients enter the analysis set (93 minus two donors lacking
the 28-day outcome label), of whom 90 have direct age records. All known ages
are adult (median 67 years, range 19--93; the two youngest donors are 19
years old). The full GSE279452 series additionally contains 26 pediatric
samples (17 sepsis, 9 healthy or ICU control; ages 0.1--15 years); none of
them entered the candidate set, so the pediatric/adult mixture and the
site--age correlation present in the full series do not enter any reported
analysis. Fatal outcomes stratify by age as expected:
0 of 6 patients below 40, 7 of 30 patients aged 40--60, and 10 of 54
patients over 60 with a known 28-day outcome (2 further donors in this stratum
lack the 28-day label; 17 fatal in total). The cohort is single-site and single-platform
(10x), which is reported as a cohort-level confound for the cross-disease
transfer analyses (Section~\ref{ssec:limitations}).

\subsection{Supplementary Note: per-patient cell-count association}
\label{sn:cellcount}

Per-patient platelet cell counts correlated only weakly with predicted scores
across models (Spearman $\rho = 0.18$--0.29 for all five foundation models and
the PCA-200 baseline; minimum 0.183 for \model{UCE}). A null model using only
cell count and disease label reached 0.554 $\pm$ 0.046 patient-level AUC within
COVID-19 (Supplementary Figure~\ref{fig:supp_cellcount}).

\subsection{Supplementary Note: platelet and megakaryocyte content of pretraining corpora}
\label{sn:corpus_platelet}

We measured or bounded the platelet content of each model's pretraining
corpus against public records. For scBaseCount (the STATE corpus component
and the full STACK corpus), the cell-type annotation space contains no
platelet or megakaryocyte category, and the QC filters ($\geq$300 genes and
$\geq$500 UMIs per cell) exclude typical anucleate platelet barcodes;
platelets are therefore structurally absent. For UCE's IMA corpus, the
published 51-category composition of the 2{,}969{,}114-cell representative
sample contains neither platelets nor megakaryocytes. For the three
Census-derived corpora, platelet-annotated cells are present and similarly
rare: approximately 0.120\% (\model{UCE}), 0.125\% (\model{scGPT}), and
0.120--0.125\% (\model{scPRINT-2}) of the respective releases, tens of
thousands of annotated cells. For \model{Geneformer} (Genecorpus-104M) and
for the full STATE corpus, cell-type-resolved composition is not published
and cannot be bounded. This measurement bears directly on the attribution
hypothesis in Section~\ref{ssec:discussion_property}: the three bounded
models share essentially identical platelet fractions yet span most of the
between-model performance spread, and the two leaders are exactly the models
that cannot be bounded, so platelet content in the corpus can neither order
the models it is measurable on nor be assessed on the models it would need
to explain.

\subsection{Supplementary Note: the 5-patient GSE167363 test set}
\label{sn:gse167363_testset}

For completeness, the cross-disease transfer analysis was also run on the
5-patient GSE167363 sepsis test set (2 fatal). Forward point estimates were
substantially higher (up to 1.00 patient-level AUC: \model{STATE} 1.00,
\model{Geneformer} 1.00, \model{scGPT} 0.833, \model{UCE} 0.667,
\model{scPRINT-2} 0.667). The reverse configuration trained on the same five
patients and tested on COVID-19 (240 patients, 25 fatal) gave \model{STATE}
0.799 (95\% CI [0.718, 0.876]) with logistic regression, \model{scPRINT-2}
0.695, \model{Geneformer} 0.668, \model{scGPT} 0.661, and \model{UCE} 0.653.
Given that the pool contains five patients, these values are reported only as
a labeled small-sample reference and must not be compared with the 91-patient
and 32-patient test sets used for all reported transfer conclusions. Two further
models, TranscriptFormer and STACK, were evaluated under a different held-out
test split; their scores are not comparable to the reported transfer and
adaptation values, and they are excluded from all such comparisons
(Section~\ref{ssec:r3_cross_disease}).

\subsection{Supplementary Note: severe-excluded sensitivity analysis}
\label{sn:severe}

Because severe cases (ICU-discharged survivors) are clinically adjacent to
fatal outcomes and severity grading is not harmonized across source cohorts,
the primary label definition includes severe cases as survivors. All headline
analyses were also computed with severe cases excluded from the survivor class;
rankings were unchanged, and absolute AUCs shifted by $<0.05$ in all cases
(Supplementary Table~\ref{tab:severe_sens}).

\subsection{Supplementary Note: aggregation sensitivity}
\label{sn:aggregation}

Patient-level scores were computed as the mean of a patient's out-of-fold cell
predictions. Recomputing all headline tables with median and 10\% trimmed-mean
aggregation changed no ranking and no significance call
(Supplementary Table~\ref{tab:aggregation_sens}).

\subsection{Supplementary Note: outcome-label provenance}
\label{sn:label_provenance}

Outcome labels for every fatal patient were verified against source records:
Ren (GEO outcome field), Stephenson (E-MTAB-10026 outcome attribute, 8 deaths),
Qiu GSE167363 (publication, 2 deaths), Liu (Table S6B, 4 deaths),
Schulte-Schrepping (Table S1, 3 deaths), Bernardes (publication, 1 death),
GSE279452 (series supplement, 17 deaths at 28 days), and Combes (Supplementary
Table 1 of the source publication). GSE151263 does not report per-patient
outcomes and is excluded from all mortality analyses
(Section~\ref{ssec:methods_cohort}).

\subsection{Supplementary Note: patient-matched monocyte comparison}
\label{sn:monocyte_matched}

The monocyte--platelet comparison was repeated on the patient intersection
identified through the verified sample-to-patient crosswalk (Ren: sample title
base to patient; Stephenson: individual identifier; GSE167363: internal sample
map). Apparent direction flips under single-seed splits were traced to
split-to-split variability and disappeared under the 10-seed ensemble:
\model{STATE}, \model{scGPT}, and the PCA-200 baseline returned to the
platelet-favoring direction (or to zero difference) in all configurations,
while \model{UCE} favored monocytes in every configuration, a model-specific
exception that is reported as such
(Section~\ref{ssec:r4_monocyte}).
